\section{What the proof actually makes available}
\label{sec:architecture}

\subsection{The endpoint and its scope}

The readable project's \src{Statements.lean} imports Mathlib and states
RH and zero-location / nonvanishing assertions for Dirichlet characters.
For example, the right-half-plane statement quantifies over every
positive modulus \(N\), every \(\C\)-valued Dirichlet character \(\chi\),
and every \(s\) with \(\Rea s>1/2\), excluding \(s=1\) for the principal
character. It does not quantify over arbitrary Hecke characters of number
fields, Dedekind zeta functions, or automorphic representations
\cite{statements}.

The repository supplies a readable and a compact edition. This report uses
the readable source names. Earlier reports sometimes quote the compact
edition's differently organized paths and slightly different intermediate
constants. Such differences are not mathematical contradictions.

\subsection{The central covolume comparison}

The arithmetic construction uses definite unitary groups over a CM field,
derivatives of theta functions, iterated theta lifts, and integral lattices
of automorphic forms. Its three central inputs are:
\begin{enumerate}[leftmargin=*,itemsep=3pt]
\item \emph{Short forms:} theta derivatives carry divisibility gains at
specified primes, while their archimedean sizes remain controlled.
\item \emph{Enough forms:} a second-moment argument supplies a nearly
spanning independent family, so a norm bound controls a substantial
covolume rather than a few isolated vectors.
\item \emph{Integral forms are long:} comparison of determinant heights,
using separately fitted twists at the coefficient primes, gives a lower
bound for the same covolume corrected by an index.
\end{enumerate}

Writing \(d\) for target dimension, \(n\) for the selected family's rank,
\(V\) for logarithmic squared covolume, \(I\) for its index bound,
and \(H_{\rm main}\) for the shared main height, the comparison has the
schematic form
\[
 d(H_{\rm main}-E_{\rm low})
 \ \le\ V+2I
 \ \le\ n(H_{\rm main}-m\,\mathrm{rem}+E_{\rm up})+2I.
\]
The main height is nonnegative, \(n\le d\), and \(n\ge d/2\).
The source's \src{star\_inequality} concludes exactly
\[
             m\,\mathrm{rem}\le f a(\log b)^5.
\]
Here \(m\) is the number of lifts, \(f\) counts infinite places, and \(a\)
is an arithmetic scale \cite{mainineq}. With \(m\asymp\sqrt{x}\),
careful conversion of the weighted remainder gives the critical
square-root saving. The readable source obtains an intermediate
\(x^{-1/2+6\varepsilon}\) bound with all parameters fixed in the
prescribed order \cite{onesided}. The exponent \(1/2\) is connected to
the chain length; a new proof cannot ignore the costs of maintaining
integrality and spanning through that chain.

This is the primary transferable idea: compare the \emph{same} integral
object from both sides, so the large common height cancels. A family
of unrelated short forms, or a determinant lower bound with incompatible
normalization, does not produce the inequality.

\subsection{The analytic endpoint}

The arithmetic coefficient in the source's prime sum is built from
\[
 \Lambda(n)\,\mathbf1_{n\equiv1\pmod q}(1-\chi_d(n)).
\]
A reciprocal-ramp kernel, \((1/n-1/x)_+\), turns its logarithmic-derivative
series into a meromorphic Mellin transform. Its poles record
differences of zero multiplicities of untwisted and quadratically twisted
Dirichlet \(L\)-functions. A one-sided bound and a Landau-type
nonnegative-transform argument exclude the relevant poles.
An auxiliary twist is chosen to avoid zeros in a specified compact
region, preventing cancellation of the zero under investigation
\cite{criterion,landau}.

The meromorphic-transform argument can be reused in broader settings.
The difficult part of a new RH proof is obtaining the new prime sum
with the correct sign and the same square-root saving. Depending on
the chosen detector, a known baseline can replace the zero-avoiding
analytic twist, but it does not replace the arithmetic inequality.

\subsection{The torsion-selection interface}

The source's \src{RootsOfUnity.lean} proves that finitely many integral
polynomials that are not excessively divisible at given points can
remain nearly as little divisible at a common tuple
\((\zeta_j-1)_j\), where the \(\zeta_j\) have \(p\)-power order.
The bound is chosen before the polynomials, and individual values incur
a positive tolerance. Its scalar step compares with a Gauss valuation
at the \emph{selected} shell radius \cite{roots}.

The arithmetic callers encode local conditions and denominator estimates
in such polynomials. This makes the selection lemma a separate
mathematical asset: its proof does not require the truth of RH.
The next section improves a different, useful comparison---with the
\emph{original fixed-radius} Gauss valuation---and proves why the distinction
is essential.

\subsection{Source map}

Paths below are relative to \texttt{grh/src/}.
The file \path{evidence/source_manifest.json} records exact URLs and hashes.

\begin{center}\small
\begin{tabularx}{\linewidth}{@{}p{0.52\linewidth}X@{}}
\toprule
Source & Role\\ \midrule
\path{Statements.lean}, \path{Main.lean} & Mathlib endpoint statements\\
\path{New/Reduction/MainInequality.lean} & Main height cancellation\\
\path{New/Reduction/ShortAndEnough.lean} & Shortness and spanning\\
\path{New/Determinant/Twist/LongForms.lean} & Covolume lower bound\\
\path{New/Reduction/OneSidedBound.lean} & Scale choice and GRH endpoint\\
\path{New/Reduction/Zeros/Criterion.lean} & Zero-multiplicity detector\\
\path{New/Congruences/Selmer/RootsOfUnity.lean} & Scalar and packet
torsion selectors\\
\path{New/Congruences/FirstOrder/DenominatorBound.lean} &
How even classes are avoided\\
\path{Lib/Cyclotomic/OddSha.lean} & Fixed-field odd-cocycle length bound\\
\bottomrule
\end{tabularx}
\end{center}

